\UseRawInputEncoding

\lstdefinestyle{promptstyle}{
  basicstyle=\ttfamily\scriptsize,
  breaklines=true,
  breakatwhitespace=false,
  columns=fullflexible,
  keepspaces=true,
  showstringspaces=false,
  tabsize=2
}

\begin{tcblisting}{
    listing only,
    breakable,
    enhanced,
    colback=lightblue!5,
    colframe=lightblue!50,
    rounded corners,
    top=1pt,bottom=1pt,left=2pt,right=2pt,boxsep=0.5pt,
    title={Prompt for Listwise Contrastive Rubric Generation ({\dataset} Curation)},
    listing options={style=promptstyle}
}
You are an expert in pedagogy and critical thinking. Your mission is to create a universal scoring rubric based on a user's request and an ordered list of example responses. The final rubric must consist of high-level, generalizable principles that can be used to evaluate any response to the request, not just the specific examples provided.

====================================================================
Methodology - A Three-Step Process for Principled Rubric Design
====================================================================

1. Step 1: Extract Explicit Requirements.  
   - Meticulously analyze the <request> tag to identify all direct commands and constraints (e.g., length, format, style).  
   - These requirements are *non-negotiable hard rules* that must appear in the rubric.  
   - They should be clearly labeled as [Hard Rule] in the final output.  

2. Step 2: Analyze the Ordered Examples for Specific Differences.  
   - Study the <responses> list, which is sorted in *descending preference* (earlier responses are BETTER).  
   - Identify concrete qualities that make higher-ranked responses superior to lower-ranked ones. Consider both adjacent comparisons (i vs. i+1) and contrasts between the top responses and the rest.  
   - It is acceptable at this stage to note topic-specific observations (e.g., "Response 1 includes citation X"), but these are *temporary* and must not appear in the final rubric.  
   - Every such observation must then be abstracted in Step 3.  

3. Step 3: MANDATORY ABSTRACTION - Convert Specifics to Universal Principles.  
   - This is the most critical step. For each observation from Step 2, ask:  
     **"What is the universal principle of high-quality communication, reasoning, or pedagogy that this specific difference demonstrates?"**  
   - Convert each observation into a principle that applies across any domain, not just the provided examples.  
   - Any rubric item that references concrete facts, names, events, topics, or response indices (e.g., "Response #1") is INVALID.  
   - All such principles must be labeled as [Principle] in the final output.  

====================================================================
Strict Guidelines for Final Output
====================================================================

- **Abstraction is Mandatory:**  
  Every rubric item must be a universal principle. If any rubric still contains topic-specific references (e.g., names, places, myths, numbers, historical facts), or mentions response indices/positions, it is automatically invalid.

- **Two Distinct Categories:**  
  - [Hard Rule]: Derived strictly from explicit requirements in the <request>.  
  - [Principle]: Derived from abstracted differences in Step 3.  

- **Comprehensiveness:**  
  The rubric must cover all critical aspects implied by the request and examples, including explicit requirements and implicit quality standards.

- **Conciseness & Uniqueness:**  
  Each rubric must capture a distinct evaluation criterion. Overlapping or redundant criteria must be merged into a single rubric. Wording must be precise and free of repetition.

- **Format Requirements:**  
  - Use a numbered list.  
  - Each item starts with "The response..." phrased in third person.  
  - Append [Hard Rule] or [Principle] at the end of each item.  
  - Do not include reasoning, explanations, or examples in the final output—only the rubrics.

- **Validation Check Before Output:**  
  Before presenting the final list, verify:  
  1. Does every rubric meet the abstraction requirement (no topic-specific details, no reference to response indices)?  
  2. Are all hard rules from Step 1 included?  
  3. Are all principles unique and non-overlapping?  
  4. Is the list written entirely in third person, concise, and consistent?

====================================================================
Final Output Format
====================================================================

1. The response ... [Hard Rule]  
2. The response ... [Principle]  
3. The response ... [Principle]  
... (continue until all rules and principles are listed)

====================================================================

<request>
{request}
</request>

<context>
{context}
</context>

<responses>
{responses}
</responses>
\end{tcblisting}
